% Werke X/1 (1917), native227–229, printed220–222.
\begingroup
\makeatletter\tagsleft@true\let\veqno\@@leqno\makeatother
\clearpage\SourceLeaf{227}\ArithmeticHead{220}{L}{ANALYSIS.\quad NACHLASS.}
\NumberedItem{4}
\HistoricTitle{Schöner Lehrsatz.}
\begin{equation*}\tag*{[2]}
\frac{\dfrac{2\partial m}m-\dfrac{2\partial n}n}{mm-nn}
=\frac{\dfrac{\pi\partial\lambda}\lambda-\dfrac{\pi\partial\mu}\mu}{\lambda\mu}.
\end{equation*}
\par\indent Integrationsfähig durch eine Reihe wird die vorstehende Gleichung, wenn man sie in folgende Form setzt:
\[\frac{2\mu\mu}{mm-nn}\left(\frac{dm}m-\frac{dn}n\right)=-\pi d\frac\mu\lambda.\]
\par\indent Das Hauptmoment des Beweises des umstehenden Lehrsatzes ist folgendes. Setzt man
\[\frac nm=a,\quad\frac{\surd(mm-nn)}m=b,\quad\frac m\lambda=A,\quad\frac m\mu=B,\]
\noindent so wird
\[\begin{aligned}
A&=1+\frac14aa+\frac14.\frac9{16}a^4+\frac{1.9.25}{4.16.64}a^6+\text{etc.}\\[.5em]
B&=1+\frac14bb+\frac14.\frac9{16}b^4+\frac{1.9.25}{4.16.64}b^6+\text{etc.}
\end{aligned}\]
\begin{equation*}\tag*{(1)}
aA+(3aa-1)\frac{dA}{da}+(a^3-a)\frac{ddA}{da^2}=0,
\end{equation*}
\[bB+(3bb-1)\frac{dB}{db}+(b^3-b)\frac{ddB}{db^2}=0,\;\text{[*)]}\]
\noindent woraus
\begin{equation*}\tag*{(2)}
aB+(3aa-1)\frac{dB}{da}+(a^3-a)\frac{ddB}{da^2}=0.
\end{equation*}
\noindent Aus (1) und (2) folgt, wenn wir
\[\frac{BdA}{da}-\frac{AdB}{da}=u\]
\noindent setzen,
\[(3aa-1)u+(a^3-a)\frac{du}{da}=0\]
\noindent oder
\[(a^3-a)u=\text{Const.}\]
\par\medskip\noindent\rule{.25\linewidth}{.4pt}
\par\smallskip\noindent [*) In der Handschrift steht \(b^3-a\) statt \(b^3-b\).]
\clearpage\SourceLeaf{228}\ArithmeticHead{221}{R}{\scriptsize DIE UNENDLICH VIELEN WERTE DES ARITHM.-GEOM.\ MITTELS.}
\par\indent Nun wird für \(a=1,\;b=0\)
\[\begin{gathered}
A=\infty,\quad(a^3-a)A=0;\\[.4em]
B=1;\quad\frac{dA}{da}=\infty,\\[.5em]
(a^3-a)\frac{dA}{da}=-\frac2\pi,\quad\frac{dB}{da}=-\frac12,\\[.5em]
(a^3-a)u=-\frac2\pi.
\end{gathered}\]
\NumberedItem{5}
\begin{center}\textls[120]{Summation einer Reihe, wo die Logarithmen der Glieder eine arithmetische Progression der zweiten Ordnung bilden.}\end{center}
\par\indent Es seyn die Logarithmen dreier auf einander folgender Glieder:
\[a-b-c,\quad a,\quad a+b-c\]
\noindent und von einer andern Reihe
\[\alpha-\beta-\gamma,\quad\alpha,\quad\alpha+\beta-\gamma.\]
\noindent Ist nun
\[\begin{aligned}
\alpha&=a+\frac{bb}{4c}+\frac12\log k\pi-\frac12\log c\\[.5em]
\beta&=i.180^0.\frac bc\\[.5em]
\gamma&=\frac{kk\pi\pi}c,
\end{aligned}\]
\[\begin{aligned}
\tfrac12\log k\pi&=0{,}0674670.920.5\\
k\pi&=1{,}3643763538.418\\
kk\pi\pi&=1{,}86152283563\\
4kk\pi\pi&=7{,}4460913425,
\end{aligned}\]
\noindent so sind die Summen beider Reihen gleich.
\clearpage\SourceLeaf{229}\ArithmeticHead{222}{L}{ANALYSIS.\quad NACHLASS.}
\begin{center}\textls[120]{Beispiel.}\end{center}
\[\begin{array}{|r|r|r||}
4{,}2&&\\
&+4{,}7&\\
8{,}9&&-3{,}6\\
&+1{,}1&\\
{}[1]0{,}0&&-3{,}6\\
&-2{,}5&\\
7{,}5&&-3{,}6\\
&-6{,}1&\\
1{,}4&&
\end{array}\]
\noindent [Es folgt eine logarithmische Rechnung für \(a=10,\;b=-0{,}7,\;c=1{,}8\), die ergibt:
\[\alpha=1{,}01832497,\quad\beta=-i.70^0,\quad\gamma=1{,}0341793.]\]
\NumberedItem{6}
\par\indent Ist
\begin{equation*}\tag*{[3]}
\frac nm=\frac{\mu(1-2e^{-M\pi}+2e^{-4M\pi}-2e^{-9M\pi}+\ldots)^2}{\mu(1+2e^{-M\pi}+2e^{-4M\pi}+2e^{-9M\pi}+\ldots)^2}
\end{equation*}
\noindent so kann man statt \(M\) setzen
\begin{equation*}\tag*{[4]}
M^*=\frac1{2ai}
\raisebox{-.8ex}{$\displaystyle+\frac1{2bi}$}
\raisebox{-1.6ex}{$\displaystyle+\frac1{2ci}$}
\raisebox{-2.4ex}{$\displaystyle+\frac1{2di}$}
\raisebox{-3.2ex}{$\displaystyle+\frac1{2ei}$}
\raisebox{-4ex}{$\displaystyle+\text{etc.}+\frac1M$}
\end{equation*}
\noindent wo \(a,b,c,d,e\) u.s.w.\ eine beliebige ungerade Menge ganzer reeller Zahlen bedeuten; oder auch
\begin{equation*}\tag*{[5]}
\frac{pM+2qi}{r+2sMi}=\frac{M+2(qr-psMM)i}{rr+4ssMM},
\end{equation*}
\noindent wo \(p,q,r,s\) beliebige der Bedingung
\[pr+4qs=1\]
\noindent genüge leistende ganze reelle Zahlen sind.
\clearpage\endgroup
